\documentclass[10pt,a4paper]{amsart}
\usepackage{amsmath,amssymb,amsthm,graphicx}
% Retained minimal preamble for the Stacks Project chapter sites-cohomology.tex (Cohomology on Sites).
% Verbatim excerpt of the chapter preamble: Xy-pic setup (native preamble lines 14-18) and the native
% theorem declarations with the mathematical macros (native preamble lines 154-217). Copied mechanically;
% nothing added, removed or redefined. The wrapper class and the amsmath/amssymb/amsthm/graphicx base
% packages are supplied by the assembler, so the class/externaldocument part of the native preamble is omitted.
% For commutative diagrams we use Xy-pic
\usepackage[all]{xy}

% We use 2cell for 2-commutative diagrams.
\xyoption{2cell}
% Theorem environments.
%
\theoremstyle{plain}
\newtheorem{theorem}[subsection]{Theorem}
\newtheorem{proposition}[subsection]{Proposition}
\newtheorem{lemma}[subsection]{Lemma}

\theoremstyle{definition}
\newtheorem{definition}[subsection]{Definition}
\newtheorem{example}[subsection]{Example}
\newtheorem{exercise}[subsection]{Exercise}
\newtheorem{situation}[subsection]{Situation}

\theoremstyle{remark}
\newtheorem{remark}[subsection]{Remark}
\newtheorem{remarks}[subsection]{Remarks}

\numberwithin{equation}{subsection}

% Macros
%
\def\lim{\mathop{\mathrm{lim}}\nolimits}
\def\colim{\mathop{\mathrm{colim}}\nolimits}
\def\Spec{\mathop{\mathrm{Spec}}}
\def\Hom{\mathop{\mathrm{Hom}}\nolimits}
\def\Ext{\mathop{\mathrm{Ext}}\nolimits}
\def\SheafHom{\mathop{\mathcal{H}\!\mathit{om}}\nolimits}
\def\SheafExt{\mathop{\mathcal{E}\!\mathit{xt}}\nolimits}
\def\Sch{\mathit{Sch}}
\def\Mor{\mathop{\mathrm{Mor}}\nolimits}
\def\Ob{\mathop{\mathrm{Ob}}\nolimits}
\def\Sh{\mathop{\mathit{Sh}}\nolimits}
\def\NL{\mathop{N\!L}\nolimits}
\def\CH{\mathop{\mathrm{CH}}\nolimits}
\def\proetale{{pro\text{-}\acute{e}tale}}
\def\etale{{\acute{e}tale}}
\def\QCoh{\mathit{QCoh}}
\def\Ker{\mathop{\mathrm{Ker}}}
\def\Im{\mathop{\mathrm{Im}}}
\def\Coker{\mathop{\mathrm{Coker}}}
\def\Coim{\mathop{\mathrm{Coim}}}

% Boxtimes
%
\DeclareMathSymbol{\boxtimes}{\mathbin}{AMSa}{"02}

%
% Macros for moduli stacks/spaces
%
\def\QCohstack{\mathcal{QC}\!\mathit{oh}}
\def\Cohstack{\mathcal{C}\!\mathit{oh}}
\def\Spacesstack{\mathcal{S}\!\mathit{paces}}
\def\Quotfunctor{\mathrm{Quot}}
\def\Hilbfunctor{\mathrm{Hilb}}
\def\Curvesstack{\mathcal{C}\!\mathit{urves}}
\def\Polarizedstack{\mathcal{P}\!\mathit{olarized}}
\def\Complexesstack{\mathcal{C}\!\mathit{omplexes}}
% \Pic is the operator that assigns to X its picard group, usage \Pic(X)
% \Picardstack_{X/B} denotes the Picard stack of X over B
% \Picardfunctor_{X/B} denotes the Picard functor of X over B
\def\Pic{\mathop{\mathrm{Pic}}\nolimits}
\def\Picardstack{\mathcal{P}\!\mathit{ic}}
\def\Picardfunctor{\mathrm{Pic}}
\def\Deformationcategory{\mathcal{D}\!\mathit{ef}}

% --- appended 2026-10-12 (Chief): the project own definitions of the reference/slogan/history
% display environments, copied mechanically from the shared Stacks preamble lines 8-12
% (\usepackage{verbatim} plus the three \newenvironment lines, which the chapter preamble keeps
% verbatim from the project preamble). Nothing added, removed or redefined: these environments
% carry no mathematical content and the project itself typesets them as comments. Without these
% three lines a counted statement whose own text uses them fails with Environment slogan undefined.
\usepackage{verbatim}
\newenvironment{reference}{\comment}{\endcomment}
\newenvironment{slogan}{\comment}{\endcomment}
\newenvironment{history}{\comment}{\endcomment}
\usepackage{cleveref}
\crefname{originaldefinition}{original definition}{original definitions}
\Crefname{originaldefinition}{Original definition}{Original definitions}
\makeatletter
\newlabel{categories-definition-pullback-functor-fibred-category}{{02XN}{source}{}{sourcepointer-1}{}}
\newlabel{categories-definition-pullback-functor-fibred-category@cref}{{[originaldefinition][0][]02XN}{[1][0][]source}{}{}{}}
\makeatother
\crefname{originallemma}{original lemma}{original lemmas}
\Crefname{originallemma}{Original lemma}{Original lemmas}
\makeatletter
\newlabel{derived-lemma-right-derived-delta-functor}{{05TE}{source}{}{sourcepointer-2}{}}
\newlabel{derived-lemma-right-derived-delta-functor@cref}{{[originallemma][0][]05TE}{[1][0][]source}{}{}{}}
\makeatother
\crefname{originalreference}{original reference}{original references}
\Crefname{originalreference}{Original reference}{Original references}
\makeatletter
\newlabel{example-left-derived-colimits}{{08PG}{source}{}{sourcepointer-3}{}}
\newlabel{example-left-derived-colimits@cref}{{[originalreference][0][]08PG}{[1][0][]source}{}{}{}}
\makeatother
\crefname{originallemma}{original lemma}{original lemmas}
\Crefname{originallemma}{Original lemma}{Original lemmas}
\makeatletter
\newlabel{homology-lemma-efface-implies-universal}{{010T}{source}{}{sourcepointer-4}{}}
\newlabel{homology-lemma-efface-implies-universal@cref}{{[originallemma][0][]010T}{[1][0][]source}{}{}{}}
\makeatother
\crefname{originallemma}{original lemma}{original lemmas}
\Crefname{originallemma}{Original lemma}{Original lemmas}
\makeatletter
\newlabel{homology-lemma-uniqueness-universal-delta-functor}{{010U}{source}{}{sourcepointer-5}{}}
\newlabel{homology-lemma-uniqueness-universal-delta-functor@cref}{{[originallemma][0][]010U}{[1][0][]source}{}{}{}}
\makeatother
\crefname{originalreference}{original reference}{original references}
\Crefname{originalreference}{Original reference}{Original references}
\makeatletter
\newlabel{lemma-compute-left-derived-pi-shriek-pre}{{08PI}{source}{}{sourcepointer-6}{}}
\newlabel{lemma-compute-left-derived-pi-shriek-pre@cref}{{[originalreference][0][]08PI}{[1][0][]source}{}{}{}}
\makeatother
\crefname{originalreference}{original reference}{original references}
\Crefname{originalreference}{Original reference}{Original references}
\makeatletter
\newlabel{lemma-compute-pi-shriek}{{08PE}{source}{}{sourcepointer-7}{}}
\newlabel{lemma-compute-pi-shriek@cref}{{[originalreference][0][]08PE}{[1][0][]source}{}{}{}}
\makeatother
\crefname{originallemma}{original lemma}{original lemmas}
\Crefname{originallemma}{Original lemma}{Original lemmas}
\makeatletter
\newlabel{sites-modules-lemma-module-quotient-flat}{{03EW}{source}{}{sourcepointer-8}{}}
\newlabel{sites-modules-lemma-module-quotient-flat@cref}{{[originallemma][0][]03EW}{[1][0][]source}{}{}{}}
\makeatother
\crefname{originalremark}{original remark}{original remarks}
\Crefname{originalremark}{Original remark}{Original remarks}
\makeatletter
\newlabel{sites-modules-remark-localize-presheaves}{{03EJ}{source}{}{sourcepointer-9}{}}
\newlabel{sites-modules-remark-localize-presheaves@cref}{{[originalremark][0][]03EJ}{[1][0][]source}{}{}{}}
\makeatother
\crefname{originalreference}{original reference}{original references}
\Crefname{originalreference}{Original reference}{Original references}
\makeatletter
\newlabel{situation-fibred-category}{{08P8}{source}{}{sourcepointer-10}{}}
\newlabel{situation-fibred-category@cref}{{[originalreference][0][]08P8}{[1][0][]source}{}{}{}}
\makeatother
\title{The higher derived lower shriek along a fibred category agrees with the explicitly constructed functors}
\author{The Stacks Project Authors}
\date{stacks-project:sites-cohomology.tex:local-snapshot-20261011; GNU FDL 1.2 or later}
\begin{document}
\maketitle
\noindent\textit{Editorial note.} All delivered mathematical text is the author's own, copied from the pinned Stacks Project chapter file sites-cohomology.tex; no statement or proof step was rewritten, reordered or removed.\par\medskip
\noindent\textit{Editorial note.} The counted proof is the author's complete argument: the proof is written out in full: the delta-functor setup, the computation of the presheaves F\_\{U'\} and of their restrictions to the fibre categories, the resulting vanishing of the higher functors, the quotient argument giving universality, and the final identification with the derived lower shriek.\par\medskip
\noindent\textit{Editorial note.} The label lemma-compute-left-derived-pi-shriek names the counted statement, and the proof environment immediately following it in the source is delivered verbatim as its proof.\par\medskip
\noindent\textit{Editorial note.} The counted statement is carried by the author's own lemma environment in the chapter source.\par\medskip
\noindent\textit{Editorial note.} Pre-open gate: python operations/agent-authoring/tools/precard\_gate.py --source sites-cohomology.tex --label lemma-compute-left-derived-pi-shriek --start 10614 --end 10675 --local-id 3548 returned RESULT=OK (L3548 free, occupancy FREE). Candidate taken from the Lead-forwarded sorted list lane-s4/halfB-priority.json.\par\medskip
\noindent\textit{Editorial note.} The statement and proof are copied from the pinned chapter file; the cited inputs are kept as external pointers with their canonical Stacks tags.\par\medskip
\noindent\textit{Allowed standard prerequisites (editorial).}\begin{itemize}
\item The explicitly constructed functors L\_n on abelian presheaves. Same chapter, the lemma constructing the functors L\_n by the cosimplicial procedure, whose agreement with the derived lower shriek is the assertion. Supplies the functors whose agreement with L\_n pi\_! is proved.
\item Vanishing of the higher homology of the localisation-by-zero presheaves. Same chapter, the example computing the higher homology of the presheaves j\_\{!\}Z over a localized site, and Modules on Sites, the lemma allowing a presheaf to be written as a quotient of one with vanishing higher L\_n. Supplies the effacing objects needed for universality.
\item Universality and uniqueness of universal delta functors. Homology, the lemma that effaceability implies universality and the lemma on uniqueness of universal delta functors; Derived Categories, the lemma comparing a right derived functor with a delta functor; Same chapter, the agreement in degree zero. Supplies the final identification of the two families of functors.
\end{itemize}\medskip
\ifdefined\hypertarget\hypertarget{sourcepointer-1}{}\fi
\noindent\textit{Source reference.} Original definition 02XN: sources/stacks-chow/tags.txt, source lines 3846--3846. Prerequisite pointer invoked at proof line(s) 10646; it names a standard fact of a sibling chapter that the counted argument uses, and every step of the counted proof built on it is retained verbatim in the delivered body.\par\medskip
\ifdefined\hypertarget\hypertarget{sourcepointer-2}{}\fi
\noindent\textit{Source reference.} Original lemma 05TE: sources/stacks-chow/tags.txt, source lines 7478--7478. Prerequisite pointer invoked at proof line(s) 10674; it names a standard fact of a sibling chapter that the counted argument uses, and every step of the counted proof built on it is retained verbatim in the delivered body.\par\medskip
\ifdefined\hypertarget\hypertarget{sourcepointer-3}{}\fi
\noindent\textit{Source reference.} Original reference 08PG: sources/stacks-sites-cohomology/pinned-original-source.tex, source lines 10044--10044. Prerequisite pointer invoked at proof line(s) 10657; it names a standard fact of the same chapter that the counted argument uses, and every step of the counted proof built on it is retained verbatim in the delivered body.\par\medskip
\ifdefined\hypertarget\hypertarget{sourcepointer-4}{}\fi
\noindent\textit{Source reference.} Original lemma 010T: sources/stacks-chow/tags.txt, source lines 1339--1339. Prerequisite pointer invoked at proof line(s) 10667; it names a standard fact of a sibling chapter that the counted argument uses, and every step of the counted proof built on it is retained verbatim in the delivered body.\par\medskip
\ifdefined\hypertarget\hypertarget{sourcepointer-5}{}\fi
\noindent\textit{Source reference.} Original lemma 010U: sources/stacks-chow/tags.txt, source lines 1340--1340. Prerequisite pointer invoked at proof line(s) 10672; it names a standard fact of a sibling chapter that the counted argument uses, and every step of the counted proof built on it is retained verbatim in the delivered body.\par\medskip
\ifdefined\hypertarget\hypertarget{sourcepointer-6}{}\fi
\noindent\textit{Source reference.} Original reference 08PI: sources/stacks-sites-cohomology/pinned-original-source.tex, source lines 10526--10526. Prerequisite pointer invoked at proof line(s) 10620; it names a standard fact of the same chapter that the counted argument uses, and every step of the counted proof built on it is retained verbatim in the delivered body.\par\medskip
\ifdefined\hypertarget\hypertarget{sourcepointer-7}{}\fi
\noindent\textit{Source reference.} Original reference 08PE: sources/stacks-sites-cohomology/pinned-original-source.tex, source lines 9960--9960. Prerequisite pointer invoked at proof line(s) 10670; it names a standard fact of the same chapter that the counted argument uses, and every step of the counted proof built on it is retained verbatim in the delivered body.\par\medskip
\ifdefined\hypertarget\hypertarget{sourcepointer-8}{}\fi
\noindent\textit{Source reference.} Original lemma 03EW: sources/stacks-chow/tags.txt, source lines 4471--4471. Prerequisite pointer invoked at proof line(s) 10662; it names a standard fact of a sibling chapter that the counted argument uses, and every step of the counted proof built on it is retained verbatim in the delivered body.\par\medskip
\ifdefined\hypertarget\hypertarget{sourcepointer-9}{}\fi
\noindent\textit{Source reference.} Original remark 03EJ: sources/stacks-chow/tags.txt, source lines 4458--4458. Prerequisite pointer invoked at proof line(s) 10636; it names a standard fact of a sibling chapter that the counted argument uses, and every step of the counted proof built on it is retained verbatim in the delivered body.\par\medskip
\ifdefined\hypertarget\hypertarget{sourcepointer-10}{}\fi
\noindent\textit{Source reference.} Original reference 08P8: sources/stacks-sites-cohomology/pinned-original-source.tex, source lines 9701--9701. Prerequisite pointer invoked at proof line(s) 10616; it names a standard fact of the same chapter that the counted argument uses, and every step of the counted proof built on it is retained verbatim in the delivered body.\par\medskip
\begin{lemma}
\label{lemma-compute-left-derived-pi-shriek}
Assumptions and notation as in Situation \ref{situation-fibred-category}.
For $\mathcal{F}$ in $\textit{Ab}(\mathcal{C})$ and $n \geq 0$
the sheaf $L_n\pi_!(\mathcal{F})$ is equal to the sheaf
$L_n(\mathcal{F})$ constructed in
Lemma \ref{lemma-compute-left-derived-pi-shriek-pre}.
\end{lemma}

\begin{proof}
Consider the sequence of functors $\mathcal{F} \mapsto L_n(\mathcal{F})$
from $\textit{PAb}(\mathcal{C}) \to \textit{Ab}(\mathcal{C})$.
Since for each $V \in \Ob(\mathcal{D})$ the sequence of functors
$H_n(\mathcal{C}_V, - )$ forms a $\delta$-functor
so do the functors $\mathcal{F} \mapsto L_n(\mathcal{F})$.
Our goal is to show these form a universal $\delta$-functor.
In order to do this we construct some abelian presheaves
on which these functors vanish.

\medskip\noindent
For $U' \in \Ob(\mathcal{C})$ consider the abelian presheaf
$\mathcal{F}_{U'} = j_{U'!}^{\textit{PAb}}\mathbf{Z}_{U'}$
(Modules on Sites, Remark \ref{sites-modules-remark-localize-presheaves}).
Recall that
$$
\mathcal{F}_{U'}(U) =
\bigoplus\nolimits_{\Mor_\mathcal{C}(U, U')} \mathbf{Z}
$$
If $U$ lies over $V = p(U)$ in $\mathcal{D})$ and $U'$ lies over $V' = p(U')$
then any morphism $a : U \to U'$ factors uniquely as $U \to h^*U' \to U'$
where $h = p(a) : V \to V'$ (see
Categories, Definition
\ref{categories-definition-pullback-functor-fibred-category}).
Hence we see that
$$
\mathcal{F}_{U'}|_{\mathcal{C}_V}
=
\bigoplus\nolimits_{h \in \Mor_\mathcal{D}(V, V')}
j_{h^*U'!}\mathbf{Z}_{h^*U'}
$$
where $j_{h^*U'} : \Sh(\mathcal{C}_V/h^*U') \to \Sh(\mathcal{C}_V)$
is the localization morphism. The sheaves $j_{h^*U'!}\mathbf{Z}_{h^*U'}$
have vanishing higher homology groups (see
Example \ref{example-left-derived-colimits}).
We conclude that $L_n(\mathcal{F}_{U'}) = 0$ for all $n > 0$ and all $U'$.
It follows that any abelian presheaf $\mathcal{F}$ is a quotient
of an abelian presheaf $\mathcal{G}$ with $L_n(\mathcal{G}) = 0$ for
all $n > 0$ (Modules on Sites, Lemma
\ref{sites-modules-lemma-module-quotient-flat}).
Since $L_n(\mathcal{F}) = L_n(\mathcal{F}^\#)$ we see
that the same thing is true for abelian sheaves. Thus
the sequence of functors $L_n(-)$ is a universal delta functor
on $\textit{Ab}(\mathcal{C})$
(Homology, Lemma \ref{homology-lemma-efface-implies-universal}).
Since we have agreement with
$H^{-n}(L\pi_!(-))$ for $n = 0$ by
Lemma \ref{lemma-compute-pi-shriek}
we conclude by uniqueness of universal $\delta$-functors
(Homology, Lemma \ref{homology-lemma-uniqueness-universal-delta-functor})
and
Derived Categories, Lemma \ref{derived-lemma-right-derived-delta-functor}.
\end{proof}

\end{document}
